\chapter{Das Kontinuum und seine Metaaussagen}

\section{Die Mächtigkeit der reellen Zahlen}

Aus den bisherigen Zahlenbänden werden die Abzählbarkeit von
\(\mathbb Q\), seine Dichtheit in \(\mathbb R\), die geordnete
Körperstruktur und die Supremumseigenschaft von \(\mathbb R\)
verwendet. Eine Theorie unendlicher Reihen wird für den folgenden
Beweis nicht vorausgesetzt.

\FormulaThmDeltaK[Die reellen Zahlen haben Mächtigkeit \(2^{\aleph_0}\)]
{|\mathbb R|=|\mathcal P(\omega)|=2^{\aleph_0}}
{B48RealContinuum}{
  \DeltaRow{Zahlbereiche}{\mathbb Q\subseteq\mathbb R}
}
\begin{tabproof}
 \proofstep{}{\BXLVIIIText{Wähle eine Aufzählung
 \((q_n)_{n<\omega}\) von \(\mathbb Q\).
 Die Abbildung
 \(r\mapsto\{n:q_n<r\}\) von \(\mathbb R\) nach
 \(\mathcal P(\omega)\) ist injektiv:
 Für \(r<s\) liegt ein rationales \(q_n\) strikt dazwischen.}}
 {\BXLVIIIWhy{Abzählbarkeit und Dichtheit der rationalen Zahlen}}
 \proofstep{}{\begin{aligned}
 s_k(S)&=\sum_{n<k}\frac{2\chi_S(n)}{3^{n+1}},
 &t(S)&=\sup_{k<\omega}s_k(S)
 \end{aligned}}
 {\BXLVIIIWhy{\(S\subseteq\omega\); endliche Summen und Supremum}}
 \proofstep{}{\BXLVIIIText{Die Partialsummen steigen und
 liegen in \([0,1]\), da die endliche geometrische Summe
 \(\sum_{n<k}2/3^{n+1}=1-3^{-k}\) ist.
 Das Supremum \(t(S)\) existiert deshalb.}}
 {\BXLVIIIWhy{Geometrische Summenformel und Vollständigkeit von \(\mathbb R\)}}
 \proofstep{4}{S\ne T,\quad m=\min(S\mathbin{\triangle}T),\quad m\in T\setminus S}
 {\rA}
 \proofstep{4}{\BXLVIIIText{Bis vor \(m\) stimmen die Summen
 überein. Die \(m\)-te Stelle liefert für \(T\) einen
 Vorsprung \(2/3^{m+1}\). Der gesamte mögliche spätere
 Beitrag von \(S\) ist höchstens \(1/3^{m+1}\), wie aus
 den endlichen geometrischen Summen durch Supremumsbildung
 folgt.}}
 {\BXLVIIIWhy{Erste unterschiedliche Ziffer und Restabschätzung}}
 \proofstep{4}{t(T)-t(S)\ge
 \frac{2}{3^{m+1}}-\frac{1}{3^{m+1}}>0}
 {\BXLVIIIWhy{Obere Schranke für \(t(S)\), untere für \(t(T)\)}}
 \proofstep{}{\mathcal P(\omega)\hookrightarrow\mathbb R,\quad
 S\longmapsto t(S)}
 {\BXLVIIIWhy{Injektivität; bei umgekehrter erster Ziffer
 \(S,T\) vertauschen}}
 \proofstep{}{|\mathbb R|=|\mathcal P(\omega)|=2^{\aleph_0}}
 {\BXLVIIIWhy{\FormulaRefAuto{B48CSB}[thm];
 \FormulaRefAuto{B48ExponentLaws}[thm]}}
\end{tabproof}
\noindent
Die ternäre Codierung mit den Ziffern \(0,2\) verhindert
die Eindeutigkeitsprobleme einer unkommentierten binären
Stellenwertdarstellung. Entscheidend ist die strikte
Restabschätzung im Beweis.

\section{Formulierungen der Kontinuumshypothese}

\FormulaDefDeltaK[Kontinuum und Kontinuumshypothese]
{\mathfrak c=|\mathbb R|=2^{\aleph_0},\qquad
 \BXLVIIICH:\Longleftrightarrow\mathfrak c=\aleph_1}
{B48CHDef}{
  \DeltaRow{Nachfolgerkardinal}{\aleph_1=\aleph_0^+}
}

\FormulaThmDeltaK[CH und das Fehlen einer Zwischenmächtigkeit]
{\BXLVIIICH\ \Longleftrightarrow\
 \neg\exists X\subseteq\mathbb R\
                  (\aleph_0<|X|<|\mathbb R|)}
{B48CHIntermediate}{
  \DeltaRow{Theorie}{\BXLVIIIZFC}
}
\begin{tabproof}
 \proofstep{1}{\BXLVIIICH}{\rA}
 \proofstep{1}{\BXLVIIIText{Zwischen \(\aleph_0\) und
 \(\aleph_1=\aleph_0^+\) liegt nach Definition keine
 Kardinalzahl. Wegen \(|\mathbb R|=\aleph_1\) kann es
 daher kein solches \(X\) geben.}}
 {\BXLVIIIWhy{Nachfolgerkardinal}}
 \proofstep{3}{\neg\BXLVIIICH}{\rA}
 \proofstep{3}{\aleph_0<\mathfrak c,\qquad
 \aleph_1\le\mathfrak c,\qquad\aleph_1<\mathfrak c}
 {\BXLVIIIWhy{Cantor, Nachfolgerkardinal und \(\mathfrak c\ne\aleph_1\)}}
 \proofstep{3}{\BXLVIIIText{Wähle eine Injektion
 \(i:\omega_1\to\mathbb R\). Ihr Bild
 \(X=i[\omega_1]\) hat Mächtigkeit \(\aleph_1\) und
 erfüllt \(\aleph_0<|X|<|\mathbb R|\).}}
 {\BXLVIIIWhy{Kardinalvergleich und Bijektion auf das Bild}}
 \proofstep{}{\BXLVIIICH\ \Longleftrightarrow\
 \neg\exists X\subseteq\mathbb R\
                 (\aleph_0<|X|<|\mathbb R|)}
 {\BXLVIIIWhy{Hinrichtung und Kontraposition der Rückrichtung}}
\end{tabproof}

\section{Gödels Richtung: CH ist nicht widerlegbar}

\FormulaThmDeltaK[Relative Konsistenz der Kontinuumshypothese]
{\BXLVIIICon(\BXLVIIIZFC)\ \Longrightarrow\
 \BXLVIIICon(\BXLVIIIZFC+\BXLVIIICH)}
{B48CHConsistent}{
  \DeltaRow{Methode}{\text{konstruktibles Universum}}
}
\begin{tabproof}
 \proofstep{1}{\BXLVIIICon(\BXLVIIIZFC)}{\rA}
 \proofstep{1}{\BXLVIIICon(\BXLVIIIZFC+\BXLVIIIGCH)}
 {\FormulaRefAuto{B48GodelConsistency}[thm]}
 \proofstep{1}{M\models\BXLVIIIZFC+\BXLVIIIGCH}
 {\FormulaRefAuto{B48Completeness}[thm]}
 \proofstep{1}{M\models 2^{\aleph_0}=\aleph_0^+=\aleph_1}
 {\BXLVIIIWhy{Instanz von GCH an \(\aleph_0\)}}
 \proofstep{1}{M\models\BXLVIIIZFC+\BXLVIIICH}
 {\FormulaRefAuto{B48CHDef}[def]}
 \proofstep{1}{\BXLVIIICon(\BXLVIIIZFC+\BXLVIIICH)}
 {\FormulaRefAuto{B48ModelConsistency}[thm]}
\end{tabproof}

\FormulaThmDeltaK[Nicht-Widerlegbarkeit von CH]
{\BXLVIIICon(\BXLVIIIZFC)\ \Longrightarrow\
 \BXLVIIIZFC\nvdash\neg\BXLVIIICH}
{B48CHNotRefutable}{
  \DeltaRow{Bedeutung}{\text{Kein ZFC-Beweis von }\neg\BXLVIIICH}
}
\begin{tabproof}
\proofstep{1}{\BXLVIIICon(\BXLVIIIZFC)}{\rA}
 \proofstep{1}{\BXLVIIICon(\BXLVIIIZFC+\BXLVIIICH)}
 {\FormulaRefAuto{B48CHConsistent}{1}}
 \proofstep{1}{\BXLVIIIZFC\nvdash\neg\BXLVIIICH}
 {\FormulaRefAuto{B48RelativeNonRefutation}{2}}
 \proofstep{}{\BXLVIIICon(\BXLVIIIZFC)\longrightarrow
   \BXLVIIIZFC\nvdash\neg\BXLVIIICH}
 {\RuleRef{RI}{\rightarrow I}\,1\!:\!3}
\end{tabproof}

\section{Cohens Richtung: CH ist nicht beweisbar}

\FormulaThmDeltaK[Relative Konsistenz der Negation von CH]
{\BXLVIIICon(\BXLVIIIZFC)\ \Longrightarrow\
 \BXLVIIICon(\BXLVIIIZFC+\neg\BXLVIIICH)}
{B48NotCHConsistent}{
  \DeltaRow{Methode}{\operatorname{Add}(\omega,\aleph_2)}
}
\begin{tabproof}
 \proofstep{1}{\BXLVIIICon(\BXLVIIIZFC)}{\rA}
 \proofstep{1}{\BXLVIIICon(\BXLVIIIZFC+\BXLVIIIGCH)}
 {\FormulaRefAuto{B48GodelConsistency}[thm]}
 \proofstep{1}{M\models\BXLVIIIZFC+\BXLVIIIGCH,\quad |M|\le\aleph_0}
 {\FormulaRefAuto{B48Completeness}[thm]}
 \proofstep{1}{G\text{ \(M\)-generisch für }
                       \operatorname{Add}(\omega,\aleph_2)^M}
 {\FormulaRefAuto{B48GenericExists}[thm]}
 \proofstep{1}{M[G]\models\BXLVIIIZFC+
                (2^{\aleph_0}=2^{\aleph_1}=\aleph_2)}
 {\FormulaRefAuto{B48PlateauModel}[thm]}
 \proofstep{1}{M[G]\models\mathfrak c=\aleph_2>\aleph_1}
 {\BXLVIIIWhy{Identifikation des Kontinuums; erhaltene Nachfolgerkardinale}}
 \proofstep{1}{M[G]\models\neg\BXLVIIICH}
 {\FormulaRefAuto{B48CHDef}[def]}
 \proofstep{1}{\BXLVIIICon(\BXLVIIIZFC+\neg\BXLVIIICH)}
 {\FormulaRefAuto{B48ModelConsistency}[thm]}
\end{tabproof}

\FormulaThmDeltaK[Nicht-Beweisbarkeit von CH]
{\BXLVIIICon(\BXLVIIIZFC)\ \Longrightarrow\
 \BXLVIIIZFC\nvdash\BXLVIIICH}
{B48CHNotProvable}{
  \DeltaRow{Bedeutung}{\text{Kein ZFC-Beweis von }\BXLVIIICH}
}
\begin{tabproof}
\proofstep{1}{\BXLVIIICon(\BXLVIIIZFC)}{\rA}
 \proofstep{1}{\BXLVIIICon(\BXLVIIIZFC+\neg\BXLVIIICH)}
 {\FormulaRefAuto{B48NotCHConsistent}{1}}
 \proofstep{1}{\BXLVIIIZFC\nvdash\BXLVIIICH}
 {\FormulaRefAuto{B48RelativeNonProof}{2}}
 \proofstep{}{\BXLVIIICon(\BXLVIIIZFC)\longrightarrow
   \BXLVIIIZFC\nvdash\BXLVIIICH}
 {\RuleRef{RI}{\rightarrow I}\,1\!:\!3}
\end{tabproof}

\section{Zusammenführung}

\FormulaThmDeltaK[Unabhängigkeit der Kontinuumshypothese von ZFC]
{\begin{gathered}
 \BXLVIIICon(\BXLVIIIZFC)\ \Longrightarrow\\
 (\BXLVIIIZFC\nvdash\BXLVIIICH)
 \ \land\
 (\BXLVIIIZFC\nvdash\neg\BXLVIIICH)
 \end{gathered}}
{B48CHIndependence}{
  \DeltaRow{Kontinuumshypothese}{\BXLVIIICH:\ 2^{\aleph_0}=\aleph_1}
}
\begin{tabproof}
\proofstep{1}{\BXLVIIICon(\BXLVIIIZFC)}{\rA}
 \proofstep{1}{\BXLVIIIZFC\nvdash\BXLVIIICH}
 {\FormulaRefAuto{B48CHNotProvable}{1}}
 \proofstep{1}{\BXLVIIIZFC\nvdash\neg\BXLVIIICH}
 {\FormulaRefAuto{B48CHNotRefutable}{1}}
 \proofstep{1}{(\BXLVIIIZFC\nvdash\BXLVIIICH)\land
   (\BXLVIIIZFC\nvdash\neg\BXLVIIICH)}
 {\RuleRef{AI}{\land I}\,2,3}
 \proofstep{}{\begin{gathered}
   \BXLVIIICon(\BXLVIIIZFC)\longrightarrow\\
   ((\BXLVIIIZFC\nvdash\BXLVIIICH)\land
    (\BXLVIIIZFC\nvdash\neg\BXLVIIICH))
   \end{gathered}}
 {\RuleRef{RI}{\rightarrow I}\,1\!:\!4}
\end{tabproof}

\noindent
Nicht-Beweisbarkeit und Nicht-Widerlegbarkeit sind zwei verschiedene
metamathematische Aussagen; ihre Beweise wurden deshalb getrennt
geführt. Der Satz behauptet nicht, CH und ihre Negation seien im
selben Modell wahr. Er besagt, dass die ZFC-Axiome unter der
Konsistenzvoraussetzung beide Möglichkeiten offenlassen.
